% Use https://ctan.org/pkg/acronym?lang=en to generate an acronym list.
% In the text, you would use:
% \ac{CPU} - the first instance of the acronym will be expanded, and all uses link back to the acronym list
% \acf{CPU} - print the expanded version of the acronym
% \Ac{CPU} or \Ac{CPU} - start the expanded version of the acronym with a capital letter.
% See the documentation linked above for details

\begin{acronym}[ABCDEFGH]
	\setlength\itemsep{.0em}
	\item \textbf{Computing}
	\acro {AI} {Artificial Intelligence}
	\acro {API} {Application Programming Interface}
	\acro {CNN} {Convolutional Neural Network}
	\acro {CPU} {Central Processing Unit}
	\acro {GPU} {graphics processing unit}
	\acro {GUI} {Graphical User Interface}
	\acro {LLM} {Large Language Model}
	\acro {ML} {Machine Learning}
	\acro {MLP} {multilayer perceptron}
	\acro {RAM} {Random Access Memory}
	\acro {RGB} {red, green, and blue}
	\vspace{1em}

	\item \textbf{Fabrication}
	\acro {ALD} {Atomic Layer Deposition}
	\acro {BOX} {burried oxide}
	\acro {DoE} {Design of Experiment}
	\acro {DWB} {Direct Wafer Bonding}
	\acro {FIB} {Focused Ion Beam}
	\acro {HSQ} {hydrogen silsesquioxane}
	\acro {ICP} {Inductively Coupled Plasma}
	\acro {MOCVD} {metalorganic chemical vapour deposition}
	\acro {PECVD} {plasma-enhanced chemical vapour deposition}
	\acro {QD} {quantum dot}
	\acro {QW} {quantum well}
	\acro {RF}  {Radio Frequency}
	\acro {RIE} {Reactive Ion Etching}
	\acro {TASE} {Template-Assisted Selective Epitaxy}
	\acro {VCSEL} {Vertical Cavity Surface Emitting Laser}
	\vspace{1em}

\end{acronym}